\section{Extreme-Weather Event Definition}\label{sec:extreme_def}

\subsection{Motivation}

The forecast-error analysis in Sections~\ref{sec:descriptive}--\ref{sec:strata} pools all hours
together, so a skill metric such as MAE reflects the average behavior of the fleet without
distinguishing ordinary conditions from operationally stressful ones. Whether HRRR forecast errors
are worse, better, or unchanged during extreme weather is a distinct question, and answering it
first requires a definition of ``extreme'' that is reproducible, transferable across the
\num{896} solar and \num{114} wind sites, and independent of the forecast error it will eventually
be compared against. A companion literature review of 50 power-systems and forecasting studies
\citep{murong2026litreview} found that meteorological rarity alone is a limited proxy for
operational stress---moderately extreme, compound events that simultaneously reduce renewable
output and raise demand can pose a greater risk to supply adequacy than any single-variable
extreme \citep{chen2026}---but that the impact-based and forecast-error-based definitions the
literature otherwise favors require data (a longitudinal loss-of-load record, or the
forecast-error history this project is still building) that are not yet available. Given that
constraint, this section adopts a meteorological percentile-threshold definition as the practical
base layer for identifying candidate extreme periods, deferring the appropriateness question to a
forecast-error comparison across the resulting regimes (Section~\ref{sec:extreme_sensitivity} and
subsequent work). The full review, including its PRISMA-ScR search strategy and the corpus behind
this summary, is available at \url{docs/Literature Review - Defining Extreme Weather.pdf}.

\subsection{Baseline dataset}

\begin{itemize}
  \item \textbf{Source:} ERA5 hourly single-level reanalysis, obtained from the Copernicus
        Climate Data Store (CDS) for a bounding box covering every solar and wind site in the
        project plus a \SI{0.5}{\degree} buffer.
  \item \textbf{Baseline period:} 1991--2020, the current \SI{30}{yr} climatological standard
        normal recommended by the World Meteorological Organization \citep{wmo2017}. ERA5's long
        record (hourly, 1940--present) supplies stable percentiles even after slicing by
        hour-of-day, which the \num{2022}--\num{2024} HRRR study period alone cannot. The
        baseline ends before the study period begins, so the reference climate used to flag
        extremes is independent of the data being classified.
  \item \textbf{Fleet:} \num{896} solar sites and \num{114} wind sites (\num{1010} total).
  \item \textbf{Variables extracted:} 2\,m temperature, 2\,m dewpoint (for relative humidity),
        surface pressure, 10\,m and 100\,m wind components ($u$, $v$), and surface solar
        radiation downwards.
\end{itemize}

\subsection{Derived variables}

Site-level hourly series are built by nearest-gridpoint selection and converted to the same
variable conventions used by the HRRR-based dataset (Section~\ref{sec:errordefs}) so that
downstream comparisons reuse identical definitions:

\begin{table}[H]
\centering
\begin{threeparttable}
\caption{Derived ERA5 site-series variables.}
\label{tab:era5_derived_vars}
\begin{tabular}{llc}
\toprule
\textbf{Column} & \textbf{Formula} & \textbf{Technology} \\
\midrule
Temperature\_C       & $t2m$ (K) $- 273.15$                         & Solar, Wind \\
RH\_pct              & Magnus formula from $t2m$, $d2m$              & Solar, Wind \\
Pressure\_mbar       & $sp$ (Pa) $/ 100$                              & Solar, Wind \\
GHI                  & $ssrd$ (\si{J/m^2}, hourly-accumulated) $/ 3600 \rightarrow$ \si{W/m^2} & Solar \\
WindSpeed\_mps       & $\sqrt{u_{10}^2 + v_{10}^2}$                   & Solar (10\,m) \\
WindSpeed\_100m\_mps & $\sqrt{u_{100}^2 + v_{100}^2}$                 & Wind (100\,m) \\
\bottomrule
\end{tabular}
\begin{tablenotes}
\scriptsize
    \item \textit{Note:} ERA5 has no native 80\,m level, so 100\,m is used directly as the
    reference level closest to the HRRR hub-height convention for the wind fleet. ERA5 has no
    direct RH field, so RH is reconstructed from temperature and dewpoint via the Magnus formula.
\end{tablenotes}
\end{threeparttable}
\end{table}

\subsection{Deseasonalization}\label{sec:deseasonalization}

A percentile computed on raw meteorology mislabels ordinary diurnal and seasonal states as
extreme---every summer afternoon becomes a ``temperature extreme,'' every night a ``low-GHI
extreme.'' \citet{brunnervoigt2024} show that applying a moving window directly to raw data mixes
the mean seasonal cycle into the extreme threshold, biasing exceedance frequency by as much as
$-75\%$ in regions with a strong seasonal gradient and weak day-to-day variability. Removing the
mean cycle first eliminates this bias and permits the wider percentile-pooling window used in
Section~\ref{sec:threshold_percentile} without penalty.

The mean cycle is the empirical average for each day-of-year $d$ and hour-of-day $h$, taken across
all $N_y$ baseline years:
\[
\bar{x}(d,h) = \frac{1}{N_y}\sum_{y=1}^{N_y} x(y,d,h), \qquad
x'(t) = x(t) - \bar{x}\big(d(t), h(t)\big).
\]
Only the mean cycle is removed; the seasonal variation in \emph{spread} is left in the anomaly
$x'$. This mean-only correction is treated as a sensitivity question and tested against a
mean-plus-variance alternative (Section~\ref{sec:extreme_sensitivity}).

\textbf{GHI} is dominated by deterministic solar geometry, so it is replaced by the clear-sky
index before deseasonalization---the ratio of observed to modeled clear-sky irradiance from the
Ineichen--Perez model \citep{ineichen2002,marchesoniacland2022}, with night-time hours (solar
zenith $\ge \SI{90}{\degree}$) excluded:
\[
k_t(t) = \frac{\text{GHI}(t)}{\text{GHI}_{\text{clear-sky}}(t)}.
\]
This isolates cloud- and aerosol-driven variability---the component relevant to solar
shortfalls---from the predictable day/night and sun-angle cycle; the mean-cycle removal above is
then applied to $k_t$ rather than to raw GHI.

\subsection{Percentile threshold and persistence criterion}\label{sec:threshold_percentile}

\textbf{Variable selection.} Of the eight available variables, three carry a clear extreme
ordering and a direct link to power-system stress and are thresholded directly: GHI (as $k_t$),
temperature, and wind speed (derived from $U10$/$V10$). Relative humidity and pressure are
reserved as compound/synoptic-context variables rather than thresholded alone, since RH saturates
at 100\% and pressure extremes are better described by rate-of-change than by level. $U10$, $V10$,
and wind direction are excluded as standalone thresholds---a signed component or a circular
bearing has no scalar extreme ordering on its own.

\textbf{Percentile levels.} Following the calendar-day percentile methodology of
\citet{zhang2005} with the window correction of \citet{brunnervoigt2024}, the threshold for a
given $(d,h)$ is the $p$-th percentile of the deseasonalized anomalies pooled from a
two-dimensional window around that calendar cell---$\pm\SI{15}{days}$ (a 31-day window, circular
across the year boundary) and $\pm\SI{1}{hr}$ (circular across the day boundary)---across all 30
baseline years:
\[
Q_p(d,h) = \operatorname{percentile}_p\big\{x'(\tau) : \tau \in \mathcal{W}(d,h)\big\}.
\]
An observation is flagged extreme when its anomaly falls beyond the relevant tail:
\[
E_{\text{high}}(t) = \mathbb{1}\big[x'(t) \ge Q_{90}(d(t),h(t))\big], \qquad
E_{\text{low}}(t) = \mathbb{1}\big[x'(t) \le Q_{10}(d(t),h(t))\big].
\]
The primary levels are the 90th/10th percentiles: this is the balance point in the heat-wave
literature between events being sufficiently ``extreme'' and sufficiently ``measurable,'' since
stricter levels combined with a persistence requirement leave too few events to analyze
\citep{perkinsalexander2013}. The 95th/5th and 99th/1st pairs are carried through as robustness
checkpoints (Section~\ref{sec:extreme_sensitivity}). Table~\ref{tab:extreme_tails} lists the
relevant tail(s) and operational rationale per variable.

\begin{table}[H]
\centering
\begin{threeparttable}
\caption{Relevant extremity tail by variable.}
\label{tab:extreme_tails}
\begin{tabular}{lll}
\toprule
\textbf{Variable} & \textbf{Tail(s) of interest} & \textbf{Operational rationale} \\
\midrule
Temperature       & Both ($\ge$90th and $\le$10th) & Heat and cold both drive load extremes and affect generation \\
Wind speed        & Low ($\le$10th)                & Low wind $\rightarrow$ generation drought / adequacy risk \\
GHI (as $k_t$)     & Low ($\le$10th)                & Low irradiance $\rightarrow$ solar-generation shortfall \\
\bottomrule
\end{tabular}
\end{threeparttable}
\end{table}

\textbf{Persistence.} Isolated single-timestep exceedances are not treated as events; the
threshold must be met for at least $D$ consecutive hours:
$\text{Event}(t) = \bigwedge_{k=0}^{D-1} E(t-k)$. The tested range is $D \in \{4,8,12,16,20,24\}$
hours (Section~\ref{sec:extreme_sensitivity}).

\subsection{Normal-reference class}\label{sec:normal_reference}

Comparing forecast error between extreme and normal conditions requires a normal-reference class
that is not itself confounded with season or background regime. An hour qualifies as
normal-reference only if it satisfies \emph{both} conditions below, following two independent
precedents from the literature:

\begin{itemize}
  \item \textbf{Climatologically central:} the anomaly falls strictly between the pair's own two
        percentile thresholds (e.g., strictly between $Q_{10}$ and $Q_{90}$ for the 90/10 pair),
        the percentile-complement approach used in machine-learning forecast evaluation against a
        deseasonalized climatology \citep{lam2023}. This condition alone already guarantees the
        hour is not itself part of any qualifying event.
  \item \textbf{Temporally local:} the hour falls within $\pm\SI{5}{days}$ of the start or end
        boundary of a qualifying extreme event at that site and variable, mirroring the
        matched-control design used in heat-wave impact studies, which compare event periods
        against the days immediately before and after each event \citep{kong2024}. This holds
        season and slow-varying background conditions approximately fixed, so a later
        extreme-versus-normal error comparison isolates the effect of extremity rather than of
        the calendar.
\end{itemize}

Normal-reference \emph{events} are counted the same way as extreme events: as distinct contiguous
spans of the resulting normal-reference mask, not as raw hours.

\subsection{Regime-conditioned forecast error (planned)}\label{sec:regime_forecast_error}

Once a site-hour is labeled extreme, normal-reference, or neither, forecast error will be
evaluated \emph{within} each regime rather than pooled, because error already grows with lead time
regardless of weather (Section~\ref{sec:leadtime}); pooling all 48 lead times would confound that
growth with the extreme-versus-normal contrast. Error metrics will therefore be computed per lead
time or in matched bands (1--6, 7--24, 25--48\,h) and compared across regimes within each band.
Because extremity is defined purely on observed meteorology, the forecast itself never enters the
partition---the independent variable (was the weather extreme?) stays distinct from the dependent
one (how large was the error?). This comparison is the object of future work; the present
section establishes only the regime definition and confirms it has adequate sample size
(Section~\ref{sec:extreme_sensitivity}).

\subsection{Limitations}\label{sec:extreme_limitations}

\begin{itemize}
  \item \textbf{Resolution mismatch.} Thresholds are derived from ERA5 ($\approx\SI{31}{km}$
        grid spacing) but will be applied to HRRR ($\approx\SI{3}{km}$), so the two distributions
        can differ, most notably for wind in complex terrain. Thresholds require bias-correction
        to the HRRR climatology per site using the \num{2022}--\num{2024} overlap before
        application; this correction has not yet been carried out.
  \item \textbf{Baseline--study offset.} The 1991--2020 baseline predates the 2022--2024 study
        period, so extreme-heat events are defined against a marginally cooler climate than the
        one being evaluated.
  \item \textbf{Mean-only deseasonalization.} Section~\ref{sec:deseasonalization} removes only the
        mean seasonal cycle, not its variance; the adequacy of this simplification is tested as a
        sensitivity variant (mean $+$ variance) rather than assumed.
  \item \textbf{Deferred comparison.} This section defines the extreme/normal regimes; it does not
        yet compare HRRR forecast error between them. That comparison, and the full
        parameter-sensitivity battery beyond the sample-size check in
        Section~\ref{sec:extreme_sensitivity}, are deferred to future work.
\end{itemize}
